% ============================================
% Response to Reviewer 1
% ============================================
\section*{Response to Reviewer 1}

We thank the reviewer for their constructive feedback. Below we address each point in detail.

% ============================================
\subsection*{R1-1. Absolute biomass as an alternative explanation for Dominance}

\reviewercomment{I found the Authors' explanation of coalescence generally convincing. However, I wonder if a simple, alternative explanation for many of the results could be that parent communities attain very different absolute densities. For example, if community A reaches 10x the overall density of community B, then even relatively rare species in community A might be initially more abundant than the dominant species in community B (because communities are combined in equal volumes, line 119). This might lead to the final community being dominated by species from A simply because species survive proportionally to their initial \emph{absolute} abundance. If I understand correctly, this scenario is different from the one tested in ED Fig.~5. High heterogeneity in absolute densities could potentially explain both the prevalence of Dominance outcomes and the difference in pairwise selection correlation between same- and cross-parent pairs. In principle, heterogeneity in absolute densities could also increase with nutrient concentration. To be clear, I have no reason to expect this pattern, but I wonder if the Authors could do more to rule it out. Since ODs were measured throughout (line 379), could those data be used to check this possibility?}

\responselabel
We thank the reviewer for raising this important alternative explanation. We tested the suggested hypothesis by asking whether parental-community OD$_{600}$ correlates with the parental dominance direction or with pairwise selection correlation.
If Dominance were primarily driven by unequal absolute biomass at mixing, the higher-OD parental community should win more frequently, and signed $\Delta\text{OD}_{A-B}$, the difference between the OD of parental community A and the OD of parental community B, should be positively associated with PDI. We observed the opposite pattern. In Base, the higher-OD parental community wins only $37\%$ of Dominance events, and the continuous relationship between signed $\Delta\text{OD}_{A-B}$ and PDI is weak. This opposite-to-prediction pattern becomes strongest in Nutr$+$, where Dominance is most frequent: the higher-OD parental community wins only $9/68 = 13.2\%$ of Dominance events ($p = 3.9 \times 10^{-10}$), and signed $\Delta\text{OD}_{A-B}$ is strongly negatively correlated with PDI ($\rho = -0.60$, $p = 4.0 \times 10^{-10}$; Response Fig.~R1-1A,B). Nutr$-$ is less informative for this directional test because both OD differences and Dominance frequency are smaller. Thus, this directional analysis argues against parental-community biomass differences as an explanation for Dominance direction.

This density-based explanation also predicts that higher-OD parental communities should show stronger within-community pairwise selection correlations. To test this, we binned parental communities into low, mid, and high OD tertiles within each medium and computed the pairwise selection correlation among species from the same parental community in each OD bin (Response Fig.~R1-1C). The cross-community correlation is not re-binned by OD and is shown as a medium-level reference. Within-community pairwise selection correlation does not show a consistent increase with parental-community OD; instead, the pattern is not consistent across media. In Nutr$+$, the slight negative OD trend mirrors the negative signed $\Delta\text{OD}_{A-B}$--PDI association in the directional analysis above, rather than the positive trend expected if higher biomass drove winning. Overall, this analysis implies that differences in parental-community biomass do not explain Dominance direction or the nutrient-dependent pattern of correlated species fates.

We added a concise note to the Results and present the primary analysis in Supplementary Note~5 and Supplementary Figs.~35--36. The Results now state: \mschange{``We next tested whether nutrient-dependent changes in parental-community properties, including biomass heterogeneity and initial pool-size effects, could explain this trend; these factors did not account for the nutrient-dependent increase in Dominance or correlated species fates (Supplementary Note~5 and Supplementary Figs.~34--36).''}

\begin{figure}[H]
\centering
\includegraphics[width=\textwidth]{Fig_R1_1A_winner_loser_OD.pdf}
\caption*{\textbf{Response Fig.~R1-1A.} \textbf{Higher parental-community OD does not preferentially predict the winner in Dominance events.} Winner OD vs loser OD scatter is shown per medium; dashed line is $y=x$. Points above the line are cases where the winning parental community had higher OD than the losing parental community. The higher-OD parental community wins only $37\%$, $37\%$, and $13\%$ of Dominance events in Nutr$-$, Base, and Nutr$+$ respectively; the strongest opposite-to-prediction pattern occurs in Nutr$+$. The displayed $p$ value is from a two-sided binomial test against the null expectation that the higher-OD parental community wins $50\%$ of Dominance events.}
\end{figure}

\begin{figure}[H]
\centering
\includegraphics[width=\textwidth]{Fig_R1_1B_OD_vs_PDI.pdf}
\caption*{\textbf{Response Fig.~R1-1B.} \textbf{Continuous OD differences do not support the prediction that the higher-OD parental community wins.} Signed $\Delta\text{OD}_{A-B}$ ($\text{OD}_{A} - \text{OD}_{B}$) vs PDI is shown per medium. PDI is the Parental Dominance Index defined in the manuscript, $\text{PDI} = (2/\pi)\arctan(u/v)$. Spearman $\rho$ is the rank correlation between signed $\Delta\text{OD}_{A-B}$ and PDI; the displayed $p$ value is from a two-sided Spearman rank-correlation test of the null hypothesis of no monotonic association. The relationship is weak in Base and strongly negative in Nutr$+$ ($\rho = -0.60$, $p = 4 \times 10^{-10}$), opposite the sign predicted if higher OD caused a parental community to win. Gray points are reflections included for symmetry.}
\end{figure}

\begin{figure}[H]
\centering
\includegraphics[width=\textwidth]{Fig_R1_1C_pairwise_corr_vs_OD.pdf}
\caption*{\textbf{Response Fig.~R1-1C.} \textbf{Higher parental-community OD does not produce a consistent increase in within-community pairwise selection correlation.} Columns show Nutr$-$, Base, and Nutr$+$, respectively. \textbf{Top row:} parental communities are grouped into low, mid, and high within-medium OD tertiles; smaller text below the group labels shows the OD$_{600}$ range of each tertile. Bars show within-community pairwise selection correlation with bootstrap 95\% CIs; the dashed gray line and shaded band show the medium-level mean cross-community correlation and its bootstrap 95\% CI. \textbf{Bottom row:} per-event parental-community observations plotted by OD and within-community pairwise selection correlation; vertical dotted lines mark the low/mid/high OD boundaries used in the top row. Spearman $\rho$ and $p$ values are shown in each panel; p-values are from two-sided Spearman rank-correlation tests of the null hypothesis of no monotonic association between OD and within-community pairwise selection correlation.}
\end{figure}

% ============================================
\subsection*{R1-2. Does pH mismatch predict Dominance?}

\reviewercomment{It seems like an even more straightforward prediction would be that Dominance should become more likely when combining parent communities with qualitatively different pH (e.g.\ alkaline and acidic) compared to cases where both parents are either acidic or alkaline. Is this the case?}

\responselabel
We thank the reviewer for suggesting this direct test. Response Fig.~R1-2 shows that Dominance frequency is not significantly different between acid--alk and same-pH parental-community pairs. For this frequency-level analysis, we compared acid--alk pairs with pooled same-pH pairs (acid--acid or alk--alk) within each medium. Nutr$-$ is excluded from this analysis because parental communities show little pH change in this condition. Compared with pooled same-pH pairs, acid--alk pairs showed a slightly higher Dominance frequency in Nutr$+$ ($79.5\%$, 35/44, versus $71.7\%$, 33/46; OR $= 1.53$, Fisher $p = 0.47$), but a slightly lower Dominance frequency in Base ($61.4\%$, 27/44, versus $69.4\%$, 25/36; OR $= 0.70$, Fisher $p = 0.49$). This result implies that pH mismatch alone is not sufficient to explain Dominance frequency: even in Nutr$+$, where Dominance is common across pH-pair types, the acid--alk enrichment is not statistically distinguishable from same-pH pairs.

\reviewercomment{To support the hypothesis that the top-down regime emerges because dominant species strongly modify the pH of the ecosystem, the Authors show that the predictive power of the parent community pH increases in the Nutr+ condition when looking at coalescence between communities where one parent is acidic and one is alkaline (lines 261, ED Fig.~8).}

Yes, considering only the pH-mismatched acid--alk group amplifies the dominance-direction signal. In Nutr$+$, the lower-pH parental community dominates 72\% of all pH-pair types (binomial $p = 3.0 \times 10^{-5}$), but the acidic parental community dominates 86.4\% of acid--alk events (38/44; binomial $p = 9.4 \times 10^{-7}$; Response Fig.~R1-2). In Base, the all-pair signal is weak (54\%, binomial $p = 0.58$), and the acid--alk-only signal remains statistically insignificant (63.6\%, 28/44; binomial $p = 0.096$). Fig.~5 and the associated predictability analysis remain based on all eligible events, including both same-pH and pH-mismatched pairs. The pH-mismatch-specific analysis is reported separately in Supplementary Note~5 and Supplementary Fig.~37.

We therefore tuned the overall claim in Results \S2.5 to integrate the pH analysis with the dominant-species predictability analysis, while making clear that pH contrast alone does not explain the full Dominance pattern. The revised Results paragraph states: \mschange{``We next asked whether environmental modification through pH could contribute to the top-down regime. In our system, dominant species are often strong pH modifiers that either acidify or alkalinize the medium (Supplementary Fig.~8). In both Base and Nutr$+$ media, when these pH-modifying species become dominant within a community, they are predictive of the community's overall pH; communities dominated by acidifiers become acidic, while those dominated by alkalizers become alkaline (Supplementary Fig.~9). In coalescence events between acidic (pH $<$ 6.5) and alkaline (pH $>$ 7.5) parental communities, the acidic community dominated in only 56\% of cases in Base medium ($n = 41$), but dominated in 91\% of cases in Nutr$+$ medium ($n = 32$, Fisher's exact test $p < 0.0001$; Extended Data Fig.~8), consistent with the hypothesis that high nutrients amplify metabolic activity, thereby intensifying pH modification and excluding pH-sensitive species. A complementary pH-pairing analysis showed that pH effects do not explain the full Dominance pattern: acid--alk pairings were not significantly enriched for Dominance relative to pairings in which both parental communities had similar pH in either Base or Nutr$+$ medium (Fisher's exact tests, $p = 0.49$ and $p = 0.47$; Supplementary Fig.~37). Together, the dominant-species and pH analyses suggest that Dominance in Nutr$+$ may reflect a top-down route to community-level selection, in which dominant taxa and associated environmental feedbacks help determine winner identity; however, pH contrast alone does not account for the full Dominance pattern, and Dominance in Base is less directly explained by the measured dominant-species and pH effects tested here.''}

\begin{figure}[H]
\centering
\includegraphics[width=\textwidth]{Fig_R1_2_acidalk_per_medium.pdf}
\caption*{\textbf{Response Fig.~R1-2.} \textbf{pH contrast alone does not significantly predict Dominance frequency, but lower-pH parental communities preferentially dominate in Nutr$+$.} \textbf{(Left two panels)} Stacked class-fraction bars (Dominance / Mixture / Restructuring) for the three pH-pair types in Base and Nutr$+$; pair types are ordered Acid--Alk, Acid--Acid, Alk--Alk, with Dominance \% labeled inside each bar. Reported class-frequency p-values are from two-sided Fisher exact tests on 2$\times$2 tables comparing Acid--Alk vs pooled same-pH (Acid--Acid or Alk--Alk) events and Dominance vs non-Dominance outcomes. (All pairs) Signed asymmetry for all pH-pair types: positive = lower-pH parental community dominated, negative = higher-pH parental community dominated. (Acid--Alk pairs only) Signed asymmetry restricted to acid--alk pairs: positive = acidic parental community dominated, negative = alkaline parental community dominated. Percentages above the signed-asymmetry panels report lower-pH or acidic parental-community dominance fractions; p-values above these panels are from two-sided binomial tests against the null expectation that the lower-pH or acidic parental community dominates 50\% of events. Dots = individual events; squares = mean $\pm$ s.e.m.}
\end{figure}

% ============================================
\subsection*{R1-3. Circularity in Fig.~5C, PDI excluding dominant species}

\reviewercomment{In Fig.~5C, it seems like there is some risk of circularity if the dominant species are included in the calculation of the PDI. Does the positive correlation under Nutr+ conditions hold up if the dominant species are excluded from this calculation?}

\responselabel
We agree that this is an important circularity check. We used two removal strategies to avoid biasing the result toward the mixed-community winner. In the first, we removed the most abundant species in the coalesced community from all three community composition vectors; in the second, stricter test, we removed the most abundant species from each parental community separately before recalculating PDI. We followed the same PDI calculation and filtering criteria as in the original Fig.~5C analysis: events are included if they are classified as Mixture or Dominance before dominant species removal.

This analysis reveals that our core conclusion remains supported: dominant-species pairwise interactions predict the community-level winning direction in Nutr$+$, but not in Base. In Base, $R^2$ drops from $0.11$ to $0.00 / 0.07$ and the slope from $0.49$ to $-0.02 / 0.33$ (Original / mix-winner removed / parental-community dominants removed). In Nutr$+$, $R^2$ drops from $0.49$ to $0.24 / 0.18$ and the slope from $1.03$ to $0.64 / 0.56$ (Response Fig.~R1-3). Using independent, non-reflected events, the Spearman correlation between dominant-species pairwise score and community-level score remains significant after both removal strategies: $\rho = 0.41$ ($p = 8.2 \times 10^{-4}$) after removing the mix-winner species and $\rho = 0.42$ ($p = 1.3 \times 10^{-3}$) after removing the dominant species from both parental communities. The corresponding linear slopes are also positive and significant ($p = 1.1 \times 10^{-4}$ and $1.5 \times 10^{-4}$). Thus, the dominant species is the primary single contributor in Nutr$+$, but the Nutr$+$ trend does not disappear when dominant species are excluded. In contrast, Base does not retain a significant correlation after removal ($p = 0.94$ and $0.12$ for the two linear slopes; Spearman $p = 1.00$ and $0.21$).

We added a concise note to Results \S2.5 and present the primary analysis in Supplementary Fig.~33. The Results now state: \mschange{``The Nutr$+$ predictability trend remained significant after recalculating PDI with dominant species excluded (Spearman rank-correlation tests: $p = 8.2 \times 10^{-4}$ after removing the coalesced-community dominant species and $p = 1.3 \times 10^{-3}$ after removing the two parental-community dominant species; Supplementary Fig.~33), indicating that it is not solely a PDI circularity artifact.''}

\begin{figure}[H]
\centering
\includegraphics[width=\textwidth]{Fig_R1_3_per_medium_scatter.pdf}
\caption*{\textbf{Response Fig.~R1-3.} \textbf{Nutr$+$ predictability persists after dominant-species removal, arguing against a PDI circularity artifact.} Per-medium scatter plots are shown with no pooling between Base and Nutr$+$. Rows: Base (MN, brown) and Nutr$+$ (HN, orange). Columns: original Fig.~5C calculation; dominant species from the coalesced community removed from all three composition vectors; and dominant species from each parental community removed before recalculating PDI. Colored points are independent coalescence events; gray points are reflected counterparts plotted for symmetry. The third column is the stricter test. $R^2$ and slope annotations are computed on the reflected plotting data, as in Fig.~5C; any displayed slope p-values are from two-sided ordinary least-squares tests of slope $=0$. Spearman $\rho$ and $p$ are computed on the independent, non-reflected events; Spearman p-values are from two-sided rank-correlation tests of the null hypothesis of no monotonic association.}
\end{figure}

% ============================================
\subsection*{R1-4. Pool-size effects}

\reviewercomment{I wondered if the Authors could provide slightly more detail on the effect of the initial pool size. Is it similarly true that there is no significant effect of initial richness in the experimental communities? By eye, it looks like Dominance might be more likely for the 24 species pool, but it is hard to tell by squinting at Fig.~1E. It might also be interesting to see the survival ratio as a function of initial richness for the experimental communities.}

\responselabel
We do not detect a significant effect of pool size on Dominance frequency in the experiment, and we thank the reviewer for prompting this more direct check. Response Fig.~R1-4C shows Dominance frequency by initial pool size within each nutrient condition; the pool-size effect is not significant either when pooled across nutrient conditions ($\chi^2 = 2.24$, $p = 0.69$) or within individual nutrient conditions (Dominance vs non-Dominance tests: Nutr$-$ $p = 0.39$, Base $p = 0.82$, Nutr$+$ $p = 0.27$). Thus, although initial pool size affects realized richness and assembly survival ratio (panels A,B), these assembly effects do not translate into a significant change in Dominance frequency.

\reviewercomment{I also wondered if the Authors could look more closely at why there is no effect in the model. Perhaps by plotting the realized richness and interaction strength in parent communities as a function of initial richness?}

To address the model part of the comment, we performed a gLV pool-size ablation study at low, intermediate, and high values of the interaction-strength parameter $\mu$, varying the initial pool size from 4 to 48 species per parental community (100 replicate assemblies per $\mu$ and pool size). The model gives a consistent interpretation: at fixed $\mu$, Dominance frequency changes only modestly with pool size compared with the strong effect of $\mu$ itself (panel E), whereas the same-parental-community versus cross-parental-community pairwise-selection gap increases with pool size (panel F). Thus, pool size can modulate the strength of correlated species fates, but it does not explain the primary interaction-strength-dependent transition in Dominance frequency.

We added this analysis to Supplementary Note 5, ``Alternative models and controls for community-level selection,'' and to Supplementary Fig.~34. The new supplementary text states that pool-size variation affects assembly richness but does not explain the Dominance-frequency pattern, and that the model pool-size ablation supports the same conclusion.

\begin{figure}[H]
\centering
\includegraphics[width=\textwidth]{pool_size_analysis.pdf}
\caption*{\textbf{Response Fig.~R1-4.} \textbf{Dominance frequency is primarily associated with nutrient condition or model interaction strength, rather than initial pool size.} \textbf{A:} Realized parental richness. \textbf{B:} ASV-based assembly survival ratio, computed as realized parental ASV richness divided by inoculated species-pool size; this ratio can exceed one because the numerator is thresholded ASV richness rather than a count restricted to the inoculated isolates. \textbf{C:} Experimental Dominance fraction by medium and initial pool size. \textbf{D:} Experimental pairwise species selection grouped by initial pool size. \textbf{E:} Model Dominance fraction grouped by $\mu$. \textbf{F:} Model pairwise species selection across pool sizes $4$--$48$.}
\end{figure}

% ============================================
\subsection*{R1-5. Similarity-metric robustness claim}

\reviewercomment{The Authors write (line 134) that ``This pattern of Dominance as the most frequent outcome was robust across variants of similarity metrics''. I liked the Author's approach to measuring outcomes, and I don't advocate for changing it, but this claim seems a little misleading. Two of four alternative similarity metrics gave different results, with Dominance being the \emph{least} likely outcome under Jensen--Shannon and Jaccard.}

\responselabel
We agree that our original robustness claim was too broad and could be misleading: the Dominance pattern is recovered by Euclidean distance and Bray--Curtis dissimilarity, but not under Jaccard index or Jensen--Shannon divergence (Extended Data Fig.~2). This pattern is consistent with the fact that the metrics emphasize different aspects of community composition: abundance-weighted parental-similarity metrics such as vector decomposition, Euclidean distance, and Bray--Curtis dissimilarity are more sensitive to quantitative dominance in relative abundance, whereas Jaccard reflects presence/absence retention and Jensen--Shannon evaluates divergence across the full abundance distribution. Thus, a coalesced community can be skewed toward one parental community in abundance space while still retaining taxa from both parental communities, or while redistributing subdominant taxa in ways that make it less skewed toward one parental community under Jaccard or Jensen--Shannon.

We revised Results \S2.1 to narrow the main-text claim to the two alternative metrics that recovered the same pattern: \mschange{``This pattern of Dominance as the most frequent outcome was robust under similarity metric choices including Euclidean distance and Bray--Curtis dissimilarity (Extended Data Fig.~2).''} We also revised the Extended Data Fig.~2 caption to explicitly discuss the divergent metrics: \mschange{``The primary Dominance pattern is recovered under similarity metric choices including Euclidean distance and Bray--Curtis dissimilarity, but not under Jensen--Shannon divergence or Jaccard index, for which the coalesced community is less often measured as skewed toward one parental community. This divergence reflects the fact that the metrics emphasize different aspects of community composition: vector decomposition, Euclidean distance, and Bray--Curtis dissimilarity capture quantitative similarity to each parental community, whereas Jaccard reflects presence/absence retention and Jensen--Shannon evaluates divergence across the full abundance distribution. Thus, Jaccard and Jensen--Shannon are best interpreted as complementary metrics rather than as failed robustness checks for the abundance-weighted Dominance classification.''}


% ============================================
\subsection*{R1-6. Gray reflected points in figures}

\reviewercomment{It took me quite a while to realize that the gray points in Figs.~1E, 4C, 5C, and 6B are simply a reflection of the colored points. Please indicate this clearly, or perhaps just remove these.}

\responselabel
We retained the reflected points and updated the captions of Figs.~1E, 4C, 5C, and 6B to state explicitly that the gray points are reflections: \mschange{``Gray points represent the symmetric counterpart of each coalescence event (community B versus A).''}

We revised the captions of Figs.~1E, 4C, 5C, and 6B as quoted above.


% ============================================
\subsection*{R1-7. Interaction matrix after assembly}

\reviewercomment{In Fig.~2A, it might be helpful to show the matrix of interaction coefficients after assembly of the parent communities (with visible block structure).}

\responselabel
We decided not to add this matrix to Fig.~2D, because Fig.~2D already focuses on pairwise selection correlation and adding interaction-matrix panels would make the main figure overly dense. Instead, we added Supplementary Fig.~27 to show the requested post-assembly interaction matrices directly. The new supplementary figure shows three representative before-versus-after assembly examples, with survivor counts annotated for each parental community, and a pooled comparison showing that assembly filters surviving species into within-community groups with lower mutual competition than the unfiltered between-community coefficients at $\mu = 0.50$.

% ============================================
\subsection*{R1-8. Fig.~2D visualization}

\reviewercomment{Fig.~2D, I was somewhat confused by the relationship between points and squares (means?). Visually, the squares do not look like they are near the means of the points. Please clarify this for readers. Additionally, the gray horizontal bars are not explained in the main text unless I missed it.}

\responselabel
The previous Fig.~2D was unclear because the individual dots and summary markers represented different display levels. In the simulation panel, the dots show a stratified subset of per-event pairwise selection correlations for visual legibility, whereas the squares and error bars show the mean $\pm$ s.e.m.\ across all simulated coalescence events. This explains why the squares need not visually coincide with the apparent center of the displayed dots. We revised the Fig.~2D caption to state this distinction explicitly and to define the gray horizontal lines as the random-selection baseline from an origin-shuffle null model: \mschange{``Individual dots show per-event pairwise selection correlations; in the simulation panel, 100 stratified events per group are displayed for legibility. Squares with error bars show mean $\pm$ s.e.m.\ across all coalescence events ($n = 1{,}200$ per group in simulations and all valid experimental events in the experimental panel). Gray horizontal lines indicate the random selection baseline from a null model in which species origin labels are shuffled (see Supplementary Information).''}


% ============================================
\subsection*{R1-9. Emphasize ``cohesion without cooperation''}

\reviewercomment{Discussion, lines 342--345: I think it is worth emphasizing how interesting and perhaps surprising it is to find ``community-level cohesion without cooperation'' as the paper by M.\ Tikhonov has it. This finding will be counterintuitive to many readers, even though the mechanism is made clear.}

\responselabel
We agree with the reviewer that this is a central and counterintuitive implication of the study. We added a dedicated short paragraph to the Discussion that frames the result in these terms, placed immediately after the paragraph reconciling community-level and species-level selection across prior studies. The new paragraph emphasizes that community-level cohesion can arise without cooperative interactions, and situates this point in relation to Tikhonov's prediction and classical random Lotka--Volterra theory.

\mschange{``Notably, and perhaps counterintuitively, our results demonstrate that community-level cohesion can emerge without cooperative interactions between species, a prediction made by Tikhonov (2016) using resource-competition models and consonant with classical random Lotka--Volterra theory (May, 1972; Grilli et al., 2017). This finding is surprising because community-level behavior is often associated with cooperative mechanisms such as cross-feeding or division of labor, whereas our results indicate that assembly history and competitive interactions can also generate correlated species fates during coalescence.''}


% ============================================
\subsection*{R1-10. Means missing in Extended Data Fig.~5C}

\reviewercomment{In ED Fig.~5C, are means missing?}

\responselabel
We thank the reviewer for pointing this out. The means are shown in Extended Data Fig.~5C as square markers with error bars, but they were not sufficiently clear in the previous version. We revised Extended Data Fig.~5 so that the mean markers and error bars are more visible, and the caption now states that individual dots show per-event correlations whereas squares with error bars show mean $\pm$ s.e.m.


% ============================================
\subsection*{R1-11. Incorrect Extended Data reference in the SI}

\reviewercomment{On p.~9 in the SI, I believe the sentence ``Simulations were performed with parental community sizes ranging from 4 to 48 species per community (Extended Data Fig.~5)'' should refer to ED Fig.~4.}

\responselabel
We thank the reviewer for catching this error and corrected the SI reference.
